% ===========================================================================
% CHAPTER 1 — INTRODUCTION  (written content, ready to use)
% ---------------------------------------------------------------------------
% This is a chapter FRAGMENT: it begins at \chapter and contains no preamble.
%   * To compile inside your thesis, % ===========================================================================
% CHAPTER 1 — INTRODUCTION  (written content, ready to use)
% ---------------------------------------------------------------------------
% This is a chapter FRAGMENT: it begins at \chapter and contains no preamble.
%   * To compile inside your thesis, % ===========================================================================
% CHAPTER 1 — INTRODUCTION  (written content, ready to use)
% ---------------------------------------------------------------------------
% This is a chapter FRAGMENT: it begins at \chapter and contains no preamble.
%   * To compile inside your thesis, % ===========================================================================
% CHAPTER 1 — INTRODUCTION  (written content, ready to use)
% ---------------------------------------------------------------------------
% This is a chapter FRAGMENT: it begins at \chapter and contains no preamble.
%   * To compile inside your thesis, \input{chapter1-introduction} (or paste
%     the body into the Introduction chapter of your template / outline).
%   * Tables use plain tabular + \hline, so no extra packages are required.
%   * A "% TODO" marks the one or two spots where you should add specifics
%     only you know (e.g. the company name, if you choose to mention it).
% ===========================================================================

\chapter{Introduction}

This chapter introduces the project reported in this thesis: a browser-based
three-dimensional (3D) defence game. During the internship the author's actual
work was Frontend development; alongside it we were exposed to, and learned
about, the discipline of Business Analysis (BA). We built this game as a single
end-to-end piece of work in which we could both apply the frontend skills
practised at the company and put the BA knowledge picked up there into real use.
This chapter explains what was built, the problem it addresses, the relationship
between the internship and this independent project, the scope of what was
delivered, and how the remainder of the report is organised.

\section{Context and Motivation}

Modern web browsers have become a capable platform for rich, interactive
software. With technologies such as WebGL and high-level libraries like
Three.js, it is now possible to render and animate real-time 3D scenes directly
in a browser tab, without installing any native application. As a result, the
role of a frontend developer increasingly extends beyond building forms,
dashboards and content pages to include real-time rendering, animation,
continuous input handling, and the performance engineering those demand.

Most introductory frontend projects, however, remain data-driven interfaces.
They demonstrate component composition and state handling, but they rarely
exercise the harder problems of interactive software: maintaining many dynamic
objects that change every frame, responding to user input with no perceptible
delay, and keeping the application smooth under load. We deliberately chose a
demanding scenario --- a real-time 3D game --- precisely because it forces these
problems into the open and therefore gives a fuller picture of what frontend
engineering involves.

A game is also an honest test of analysis, not only of implementation. Before a
single line of rendering code is written, someone has to decide how the game
behaves: how enemies move, how damage is calculated, when the player wins or
loses, and how players are ranked against one another. These are precisely the
questions a Business Analyst answers --- eliciting needs, defining rules, and
writing them down so unambiguously that they can be implemented and tested.
Building a game from scratch thus let us do two things at once: practise the
frontend engineering that was the author's day-to-day internship work, and
exercise first-hand the business-analysis thinking we had learned about there
but had not performed as a formal role --- specifying the system as an analyst,
and then constructing it as a frontend engineer.

\section{Problem Statement}

The challenge this project addresses is twofold, mirroring the two skills it
sets out to demonstrate.

The \emph{analysis} problem is to take a vague idea --- ``a game where you
defend a house from monsters'' --- and turn it into a complete, consistent and
testable specification. This means identifying the users, decomposing the
system into subsystems, and pinning down every rule (how much health an enemy
has, how damage falls off with distance, how the leaderboard orders players)
with enough precision that the rules can be verified rather than argued about.

The \emph{engineering} problem is to realise that specification as a working
browser application that processes mouse input in real time, simulates and
renders dozens of moving entities (enemies, projectiles, effects) every frame,
provides immediate and legible visual feedback, and remains responsive
throughout a play session. The application must also persist results so that
players can be ranked, which introduces a small but real server-side component.

Solving both problems within one project, and keeping them traceable to each
other, is the core task of this thesis.

\section{Internship Background and Project Framing}

It is important to be clear about how this project relates to the internship.
During the internship the author's actual role and day-to-day work was
\emph{Frontend development}; that is where we gained hands-on experience and
built working habits. Business Analysis was not part of the assigned duties, but
we were exposed to it in that environment and took the opportunity to
\emph{learn about} how analysts elicit needs and specify systems. The work done
at the company is separate from this report.
% TODO (optional): if you wish to name the company or briefly describe the
% kind of frontend work done there for context, add one or two sentences.

The 3D House Defense Game presented in this thesis is therefore a
\emph{self-initiated, independent project}. It is not a company deliverable and
was not assigned by an employer; we designed and built it on our own. It served
two purposes at once: a place to apply the frontend skills practised at the
company, and --- just as valuable --- a chance to perform Business Analysis
directly, turning knowledge we had only observed during the internship into a
real, hands-on deliverable. In doing so we deliberately took on both roles:

\begin{itemize}
  \item As the \textbf{Business Analyst}, we identified the users, designed the
        gameplay and its rules, defined the leaderboard ranking policy, and
        captured everything in a formal Software Requirements Specification
        (SRS) with acceptance criteria for each requirement.
  \item As the \textbf{Frontend engineer}, we turned that SRS into a working
        application built with Next.js, React, TypeScript and Three.js, backed
        by a small persistence service for the leaderboard.
\end{itemize}

Carrying both roles was itself instructive: writing the specification first
made the implementation more disciplined, while implementing the specification
exposed gaps and over-constraints that fed back into the design. Table~1.1
summarises how the internship experience --- frontend work performed, and BA
learned about --- maps onto this independent project.

\begin{table}[h]
\centering
\caption{Internship experience applied to this independent project}
\begin{tabular}{|p{3.6cm}|p{5.0cm}|p{4.0cm}|}
\hline
\textbf{Skill area (internship exposure)} &
\textbf{How it is applied in this project} &
\textbf{Resulting deliverable} \\
\hline
Frontend development \newline (the internship work) &
Component-based UI, centralised state management, a real-time render loop, WebGL
rendering with Three.js, and a small HTTP/database service &
The implemented 3D game client and the leaderboard service (Chapter~4) \\
\hline
Business Analysis \newline (learned about, then practised here) &
Eliciting users and needs; decomposing the game into subsystems; writing
functional requirements, business rules and acceptance criteria; designing the
ranking policy &
The Software Requirements Specification and the game design (Chapter~3;
Appendix~A) \\
\hline
\end{tabular}
\end{table}

\section{Scope and Contributions}

This thesis covers a single, complete piece of work that we produced
independently. Its concrete contributions are:

\begin{itemize}
  \item \textbf{A complete Software Requirements Specification.} A structured
        catalogue of functional requirements, business rules, timing rules and
        non-functional requirements, each with acceptance criteria, covering
        every subsystem of the game (Appendix~A).
  \item \textbf{A game design derived from that specification.} The gameplay
        loop, the enemy and weapon balance, the area-of-effect damage model, the
        three-wave progression, and the leaderboard ranking rules (Chapter~3).
  \item \textbf{A working frontend application.} A real-time, interactive 3D
        game implemented with Next.js, React, TypeScript and Three.js/WebGL,
        organised into a clear layered architecture (Chapter~4).
  \item \textbf{A persistent leaderboard service.} A small server-side component
        --- HTTP endpoints backed by an SQLite database --- that records
        finished runs and returns them in ranked order (Chapter~4).
\end{itemize}

The scope is intentionally bounded. The game is a \emph{single-player},
\emph{desktop-browser} experience of \emph{three} waves, controlled with the
mouse. The following are explicitly out of scope: multiplayer or networked play,
user accounts and authentication, cross-device synchronisation, and mobile or
touch support. These boundaries keep the project achievable while still
exercising the full BA-to-implementation path, and several of them are revisited
as future work in Chapter~6.

\section{Report Structure}

The remainder of this report is organised as follows.

\begin{itemize}
  \item \textbf{Chapter~2 (Objectives)} states the general aim, the specific BA
        and Frontend objectives, and the success criteria used to judge the
        result.
  \item \textbf{Chapter~3 (Game Design)} presents the analysis work: the game
        concept, the requirements-engineering approach, the gameplay and
        economy design, the user-experience design, and the leaderboard ranking
        rules. It is the narrative companion to the full SRS in Appendix~A.
  \item \textbf{Chapter~4 (Materials and Methods)} describes the implementation:
        the technology stack and why each tool was chosen, the layered
        architecture and data flow, the state-management and game-loop design,
        the gameplay systems, the Three.js rendering layer, and the leaderboard
        service.
  \item \textbf{Chapter~5 (Results and Discussion)} evaluates the finished
        system against the specification through a requirements-traceability
        matrix, a gameplay walkthrough, and a performance assessment, and then
        reflects on the BA and Frontend decisions, the challenges met, and the
        project's limitations.
  \item \textbf{Chapter~6 (Conclusion and Perspective)} summarises what was
        achieved, names the BA and Frontend skills the project demonstrates, and
        outlines directions for future work.
  \item The \textbf{appendices} contain the full SRS, selected source-code
        listings, and a reference for the leaderboard API and database schema.
\end{itemize}
 (or paste
%     the body into the Introduction chapter of your template / outline).
%   * Tables use plain tabular + \hline, so no extra packages are required.
%   * A "% TODO" marks the one or two spots where you should add specifics
%     only you know (e.g. the company name, if you choose to mention it).
% ===========================================================================

\chapter{Introduction}

This chapter introduces the project reported in this thesis: a browser-based
three-dimensional (3D) defence game. During the internship the author's actual
work was Frontend development; alongside it we were exposed to, and learned
about, the discipline of Business Analysis (BA). We built this game as a single
end-to-end piece of work in which we could both apply the frontend skills
practised at the company and put the BA knowledge picked up there into real use.
This chapter explains what was built, the problem it addresses, the relationship
between the internship and this independent project, the scope of what was
delivered, and how the remainder of the report is organised.

\section{Context and Motivation}

Modern web browsers have become a capable platform for rich, interactive
software. With technologies such as WebGL and high-level libraries like
Three.js, it is now possible to render and animate real-time 3D scenes directly
in a browser tab, without installing any native application. As a result, the
role of a frontend developer increasingly extends beyond building forms,
dashboards and content pages to include real-time rendering, animation,
continuous input handling, and the performance engineering those demand.

Most introductory frontend projects, however, remain data-driven interfaces.
They demonstrate component composition and state handling, but they rarely
exercise the harder problems of interactive software: maintaining many dynamic
objects that change every frame, responding to user input with no perceptible
delay, and keeping the application smooth under load. We deliberately chose a
demanding scenario --- a real-time 3D game --- precisely because it forces these
problems into the open and therefore gives a fuller picture of what frontend
engineering involves.

A game is also an honest test of analysis, not only of implementation. Before a
single line of rendering code is written, someone has to decide how the game
behaves: how enemies move, how damage is calculated, when the player wins or
loses, and how players are ranked against one another. These are precisely the
questions a Business Analyst answers --- eliciting needs, defining rules, and
writing them down so unambiguously that they can be implemented and tested.
Building a game from scratch thus let us do two things at once: practise the
frontend engineering that was the author's day-to-day internship work, and
exercise first-hand the business-analysis thinking we had learned about there
but had not performed as a formal role --- specifying the system as an analyst,
and then constructing it as a frontend engineer.

\section{Problem Statement}

The challenge this project addresses is twofold, mirroring the two skills it
sets out to demonstrate.

The \emph{analysis} problem is to take a vague idea --- ``a game where you
defend a house from monsters'' --- and turn it into a complete, consistent and
testable specification. This means identifying the users, decomposing the
system into subsystems, and pinning down every rule (how much health an enemy
has, how damage falls off with distance, how the leaderboard orders players)
with enough precision that the rules can be verified rather than argued about.

The \emph{engineering} problem is to realise that specification as a working
browser application that processes mouse input in real time, simulates and
renders dozens of moving entities (enemies, projectiles, effects) every frame,
provides immediate and legible visual feedback, and remains responsive
throughout a play session. The application must also persist results so that
players can be ranked, which introduces a small but real server-side component.

Solving both problems within one project, and keeping them traceable to each
other, is the core task of this thesis.

\section{Internship Background and Project Framing}

It is important to be clear about how this project relates to the internship.
During the internship the author's actual role and day-to-day work was
\emph{Frontend development}; that is where we gained hands-on experience and
built working habits. Business Analysis was not part of the assigned duties, but
we were exposed to it in that environment and took the opportunity to
\emph{learn about} how analysts elicit needs and specify systems. The work done
at the company is separate from this report.
% TODO (optional): if you wish to name the company or briefly describe the
% kind of frontend work done there for context, add one or two sentences.

The 3D House Defense Game presented in this thesis is therefore a
\emph{self-initiated, independent project}. It is not a company deliverable and
was not assigned by an employer; we designed and built it on our own. It served
two purposes at once: a place to apply the frontend skills practised at the
company, and --- just as valuable --- a chance to perform Business Analysis
directly, turning knowledge we had only observed during the internship into a
real, hands-on deliverable. In doing so we deliberately took on both roles:

\begin{itemize}
  \item As the \textbf{Business Analyst}, we identified the users, designed the
        gameplay and its rules, defined the leaderboard ranking policy, and
        captured everything in a formal Software Requirements Specification
        (SRS) with acceptance criteria for each requirement.
  \item As the \textbf{Frontend engineer}, we turned that SRS into a working
        application built with Next.js, React, TypeScript and Three.js, backed
        by a small persistence service for the leaderboard.
\end{itemize}

Carrying both roles was itself instructive: writing the specification first
made the implementation more disciplined, while implementing the specification
exposed gaps and over-constraints that fed back into the design. Table~1.1
summarises how the internship experience --- frontend work performed, and BA
learned about --- maps onto this independent project.

\begin{table}[h]
\centering
\caption{Internship experience applied to this independent project}
\begin{tabular}{|p{3.6cm}|p{5.0cm}|p{4.0cm}|}
\hline
\textbf{Skill area (internship exposure)} &
\textbf{How it is applied in this project} &
\textbf{Resulting deliverable} \\
\hline
Frontend development \newline (the internship work) &
Component-based UI, centralised state management, a real-time render loop, WebGL
rendering with Three.js, and a small HTTP/database service &
The implemented 3D game client and the leaderboard service (Chapter~4) \\
\hline
Business Analysis \newline (learned about, then practised here) &
Eliciting users and needs; decomposing the game into subsystems; writing
functional requirements, business rules and acceptance criteria; designing the
ranking policy &
The Software Requirements Specification and the game design (Chapter~3;
Appendix~A) \\
\hline
\end{tabular}
\end{table}

\section{Scope and Contributions}

This thesis covers a single, complete piece of work that we produced
independently. Its concrete contributions are:

\begin{itemize}
  \item \textbf{A complete Software Requirements Specification.} A structured
        catalogue of functional requirements, business rules, timing rules and
        non-functional requirements, each with acceptance criteria, covering
        every subsystem of the game (Appendix~A).
  \item \textbf{A game design derived from that specification.} The gameplay
        loop, the enemy and weapon balance, the area-of-effect damage model, the
        three-wave progression, and the leaderboard ranking rules (Chapter~3).
  \item \textbf{A working frontend application.} A real-time, interactive 3D
        game implemented with Next.js, React, TypeScript and Three.js/WebGL,
        organised into a clear layered architecture (Chapter~4).
  \item \textbf{A persistent leaderboard service.} A small server-side component
        --- HTTP endpoints backed by an SQLite database --- that records
        finished runs and returns them in ranked order (Chapter~4).
\end{itemize}

The scope is intentionally bounded. The game is a \emph{single-player},
\emph{desktop-browser} experience of \emph{three} waves, controlled with the
mouse. The following are explicitly out of scope: multiplayer or networked play,
user accounts and authentication, cross-device synchronisation, and mobile or
touch support. These boundaries keep the project achievable while still
exercising the full BA-to-implementation path, and several of them are revisited
as future work in Chapter~6.

\section{Report Structure}

The remainder of this report is organised as follows.

\begin{itemize}
  \item \textbf{Chapter~2 (Objectives)} states the general aim, the specific BA
        and Frontend objectives, and the success criteria used to judge the
        result.
  \item \textbf{Chapter~3 (Game Design)} presents the analysis work: the game
        concept, the requirements-engineering approach, the gameplay and
        economy design, the user-experience design, and the leaderboard ranking
        rules. It is the narrative companion to the full SRS in Appendix~A.
  \item \textbf{Chapter~4 (Materials and Methods)} describes the implementation:
        the technology stack and why each tool was chosen, the layered
        architecture and data flow, the state-management and game-loop design,
        the gameplay systems, the Three.js rendering layer, and the leaderboard
        service.
  \item \textbf{Chapter~5 (Results and Discussion)} evaluates the finished
        system against the specification through a requirements-traceability
        matrix, a gameplay walkthrough, and a performance assessment, and then
        reflects on the BA and Frontend decisions, the challenges met, and the
        project's limitations.
  \item \textbf{Chapter~6 (Conclusion and Perspective)} summarises what was
        achieved, names the BA and Frontend skills the project demonstrates, and
        outlines directions for future work.
  \item The \textbf{appendices} contain the full SRS, selected source-code
        listings, and a reference for the leaderboard API and database schema.
\end{itemize}
 (or paste
%     the body into the Introduction chapter of your template / outline).
%   * Tables use plain tabular + \hline, so no extra packages are required.
%   * A "% TODO" marks the one or two spots where you should add specifics
%     only you know (e.g. the company name, if you choose to mention it).
% ===========================================================================

\chapter{Introduction}

This chapter introduces the project reported in this thesis: a browser-based
three-dimensional (3D) defence game. During the internship the author's actual
work was Frontend development; alongside it we were exposed to, and learned
about, the discipline of Business Analysis (BA). We built this game as a single
end-to-end piece of work in which we could both apply the frontend skills
practised at the company and put the BA knowledge picked up there into real use.
This chapter explains what was built, the problem it addresses, the relationship
between the internship and this independent project, the scope of what was
delivered, and how the remainder of the report is organised.

\section{Context and Motivation}

Modern web browsers have become a capable platform for rich, interactive
software. With technologies such as WebGL and high-level libraries like
Three.js, it is now possible to render and animate real-time 3D scenes directly
in a browser tab, without installing any native application. As a result, the
role of a frontend developer increasingly extends beyond building forms,
dashboards and content pages to include real-time rendering, animation,
continuous input handling, and the performance engineering those demand.

Most introductory frontend projects, however, remain data-driven interfaces.
They demonstrate component composition and state handling, but they rarely
exercise the harder problems of interactive software: maintaining many dynamic
objects that change every frame, responding to user input with no perceptible
delay, and keeping the application smooth under load. We deliberately chose a
demanding scenario --- a real-time 3D game --- precisely because it forces these
problems into the open and therefore gives a fuller picture of what frontend
engineering involves.

A game is also an honest test of analysis, not only of implementation. Before a
single line of rendering code is written, someone has to decide how the game
behaves: how enemies move, how damage is calculated, when the player wins or
loses, and how players are ranked against one another. These are precisely the
questions a Business Analyst answers --- eliciting needs, defining rules, and
writing them down so unambiguously that they can be implemented and tested.
Building a game from scratch thus let us do two things at once: practise the
frontend engineering that was the author's day-to-day internship work, and
exercise first-hand the business-analysis thinking we had learned about there
but had not performed as a formal role --- specifying the system as an analyst,
and then constructing it as a frontend engineer.

\section{Problem Statement}

The challenge this project addresses is twofold, mirroring the two skills it
sets out to demonstrate.

The \emph{analysis} problem is to take a vague idea --- ``a game where you
defend a house from monsters'' --- and turn it into a complete, consistent and
testable specification. This means identifying the users, decomposing the
system into subsystems, and pinning down every rule (how much health an enemy
has, how damage falls off with distance, how the leaderboard orders players)
with enough precision that the rules can be verified rather than argued about.

The \emph{engineering} problem is to realise that specification as a working
browser application that processes mouse input in real time, simulates and
renders dozens of moving entities (enemies, projectiles, effects) every frame,
provides immediate and legible visual feedback, and remains responsive
throughout a play session. The application must also persist results so that
players can be ranked, which introduces a small but real server-side component.

Solving both problems within one project, and keeping them traceable to each
other, is the core task of this thesis.

\section{Internship Background and Project Framing}

It is important to be clear about how this project relates to the internship.
During the internship the author's actual role and day-to-day work was
\emph{Frontend development}; that is where we gained hands-on experience and
built working habits. Business Analysis was not part of the assigned duties, but
we were exposed to it in that environment and took the opportunity to
\emph{learn about} how analysts elicit needs and specify systems. The work done
at the company is separate from this report.
% TODO (optional): if you wish to name the company or briefly describe the
% kind of frontend work done there for context, add one or two sentences.

The 3D House Defense Game presented in this thesis is therefore a
\emph{self-initiated, independent project}. It is not a company deliverable and
was not assigned by an employer; we designed and built it on our own. It served
two purposes at once: a place to apply the frontend skills practised at the
company, and --- just as valuable --- a chance to perform Business Analysis
directly, turning knowledge we had only observed during the internship into a
real, hands-on deliverable. In doing so we deliberately took on both roles:

\begin{itemize}
  \item As the \textbf{Business Analyst}, we identified the users, designed the
        gameplay and its rules, defined the leaderboard ranking policy, and
        captured everything in a formal Software Requirements Specification
        (SRS) with acceptance criteria for each requirement.
  \item As the \textbf{Frontend engineer}, we turned that SRS into a working
        application built with Next.js, React, TypeScript and Three.js, backed
        by a small persistence service for the leaderboard.
\end{itemize}

Carrying both roles was itself instructive: writing the specification first
made the implementation more disciplined, while implementing the specification
exposed gaps and over-constraints that fed back into the design. Table~1.1
summarises how the internship experience --- frontend work performed, and BA
learned about --- maps onto this independent project.

\begin{table}[h]
\centering
\caption{Internship experience applied to this independent project}
\begin{tabular}{|p{3.6cm}|p{5.0cm}|p{4.0cm}|}
\hline
\textbf{Skill area (internship exposure)} &
\textbf{How it is applied in this project} &
\textbf{Resulting deliverable} \\
\hline
Frontend development \newline (the internship work) &
Component-based UI, centralised state management, a real-time render loop, WebGL
rendering with Three.js, and a small HTTP/database service &
The implemented 3D game client and the leaderboard service (Chapter~4) \\
\hline
Business Analysis \newline (learned about, then practised here) &
Eliciting users and needs; decomposing the game into subsystems; writing
functional requirements, business rules and acceptance criteria; designing the
ranking policy &
The Software Requirements Specification and the game design (Chapter~3;
Appendix~A) \\
\hline
\end{tabular}
\end{table}

\section{Scope and Contributions}

This thesis covers a single, complete piece of work that we produced
independently. Its concrete contributions are:

\begin{itemize}
  \item \textbf{A complete Software Requirements Specification.} A structured
        catalogue of functional requirements, business rules, timing rules and
        non-functional requirements, each with acceptance criteria, covering
        every subsystem of the game (Appendix~A).
  \item \textbf{A game design derived from that specification.} The gameplay
        loop, the enemy and weapon balance, the area-of-effect damage model, the
        three-wave progression, and the leaderboard ranking rules (Chapter~3).
  \item \textbf{A working frontend application.} A real-time, interactive 3D
        game implemented with Next.js, React, TypeScript and Three.js/WebGL,
        organised into a clear layered architecture (Chapter~4).
  \item \textbf{A persistent leaderboard service.} A small server-side component
        --- HTTP endpoints backed by an SQLite database --- that records
        finished runs and returns them in ranked order (Chapter~4).
\end{itemize}

The scope is intentionally bounded. The game is a \emph{single-player},
\emph{desktop-browser} experience of \emph{three} waves, controlled with the
mouse. The following are explicitly out of scope: multiplayer or networked play,
user accounts and authentication, cross-device synchronisation, and mobile or
touch support. These boundaries keep the project achievable while still
exercising the full BA-to-implementation path, and several of them are revisited
as future work in Chapter~6.

\section{Report Structure}

The remainder of this report is organised as follows.

\begin{itemize}
  \item \textbf{Chapter~2 (Objectives)} states the general aim, the specific BA
        and Frontend objectives, and the success criteria used to judge the
        result.
  \item \textbf{Chapter~3 (Game Design)} presents the analysis work: the game
        concept, the requirements-engineering approach, the gameplay and
        economy design, the user-experience design, and the leaderboard ranking
        rules. It is the narrative companion to the full SRS in Appendix~A.
  \item \textbf{Chapter~4 (Materials and Methods)} describes the implementation:
        the technology stack and why each tool was chosen, the layered
        architecture and data flow, the state-management and game-loop design,
        the gameplay systems, the Three.js rendering layer, and the leaderboard
        service.
  \item \textbf{Chapter~5 (Results and Discussion)} evaluates the finished
        system against the specification through a requirements-traceability
        matrix, a gameplay walkthrough, and a performance assessment, and then
        reflects on the BA and Frontend decisions, the challenges met, and the
        project's limitations.
  \item \textbf{Chapter~6 (Conclusion and Perspective)} summarises what was
        achieved, names the BA and Frontend skills the project demonstrates, and
        outlines directions for future work.
  \item The \textbf{appendices} contain the full SRS, selected source-code
        listings, and a reference for the leaderboard API and database schema.
\end{itemize}
 (or paste
%     the body into the Introduction chapter of your template / outline).
%   * Tables use plain tabular + \hline, so no extra packages are required.
%   * A "% TODO" marks the one or two spots where you should add specifics
%     only you know (e.g. the company name, if you choose to mention it).
% ===========================================================================

\chapter{Introduction}

This chapter introduces the project reported in this thesis: a browser-based
three-dimensional (3D) defence game. During the internship the author's actual
work was Frontend development; alongside it we were exposed to, and learned
about, the discipline of Business Analysis (BA). We built this game as a single
end-to-end piece of work in which we could both apply the frontend skills
practised at the company and put the BA knowledge picked up there into real use.
This chapter explains what was built, the problem it addresses, the relationship
between the internship and this independent project, the scope of what was
delivered, and how the remainder of the report is organised.

\section{Context and Motivation}

Modern web browsers have become a capable platform for rich, interactive
software. With technologies such as WebGL and high-level libraries like
Three.js, it is now possible to render and animate real-time 3D scenes directly
in a browser tab, without installing any native application. As a result, the
role of a frontend developer increasingly extends beyond building forms,
dashboards and content pages to include real-time rendering, animation,
continuous input handling, and the performance engineering those demand.

Most introductory frontend projects, however, remain data-driven interfaces.
They demonstrate component composition and state handling, but they rarely
exercise the harder problems of interactive software: maintaining many dynamic
objects that change every frame, responding to user input with no perceptible
delay, and keeping the application smooth under load. We deliberately chose a
demanding scenario --- a real-time 3D game --- precisely because it forces these
problems into the open and therefore gives a fuller picture of what frontend
engineering involves.

A game is also an honest test of analysis, not only of implementation. Before a
single line of rendering code is written, someone has to decide how the game
behaves: how enemies move, how damage is calculated, when the player wins or
loses, and how players are ranked against one another. These are precisely the
questions a Business Analyst answers --- eliciting needs, defining rules, and
writing them down so unambiguously that they can be implemented and tested.
Building a game from scratch thus let us do two things at once: practise the
frontend engineering that was the author's day-to-day internship work, and
exercise first-hand the business-analysis thinking we had learned about there
but had not performed as a formal role --- specifying the system as an analyst,
and then constructing it as a frontend engineer.

\section{Problem Statement}

The challenge this project addresses is twofold, mirroring the two skills it
sets out to demonstrate.

The \emph{analysis} problem is to take a vague idea --- ``a game where you
defend a house from monsters'' --- and turn it into a complete, consistent and
testable specification. This means identifying the users, decomposing the
system into subsystems, and pinning down every rule (how much health an enemy
has, how damage falls off with distance, how the leaderboard orders players)
with enough precision that the rules can be verified rather than argued about.

The \emph{engineering} problem is to realise that specification as a working
browser application that processes mouse input in real time, simulates and
renders dozens of moving entities (enemies, projectiles, effects) every frame,
provides immediate and legible visual feedback, and remains responsive
throughout a play session. The application must also persist results so that
players can be ranked, which introduces a small but real server-side component.

Solving both problems within one project, and keeping them traceable to each
other, is the core task of this thesis.

\section{Internship Background and Project Framing}

It is important to be clear about how this project relates to the internship.
During the internship the author's actual role and day-to-day work was
\emph{Frontend development}; that is where we gained hands-on experience and
built working habits. Business Analysis was not part of the assigned duties, but
we were exposed to it in that environment and took the opportunity to
\emph{learn about} how analysts elicit needs and specify systems. The work done
at the company is separate from this report.
% TODO (optional): if you wish to name the company or briefly describe the
% kind of frontend work done there for context, add one or two sentences.

The 3D House Defense Game presented in this thesis is therefore a
\emph{self-initiated, independent project}. It is not a company deliverable and
was not assigned by an employer; we designed and built it on our own. It served
two purposes at once: a place to apply the frontend skills practised at the
company, and --- just as valuable --- a chance to perform Business Analysis
directly, turning knowledge we had only observed during the internship into a
real, hands-on deliverable. In doing so we deliberately took on both roles:

\begin{itemize}
  \item As the \textbf{Business Analyst}, we identified the users, designed the
        gameplay and its rules, defined the leaderboard ranking policy, and
        captured everything in a formal Software Requirements Specification
        (SRS) with acceptance criteria for each requirement.
  \item As the \textbf{Frontend engineer}, we turned that SRS into a working
        application built with Next.js, React, TypeScript and Three.js, backed
        by a small persistence service for the leaderboard.
\end{itemize}

Carrying both roles was itself instructive: writing the specification first
made the implementation more disciplined, while implementing the specification
exposed gaps and over-constraints that fed back into the design. Table~1.1
summarises how the internship experience --- frontend work performed, and BA
learned about --- maps onto this independent project.

\begin{table}[h]
\centering
\caption{Internship experience applied to this independent project}
\begin{tabular}{|p{3.6cm}|p{5.0cm}|p{4.0cm}|}
\hline
\textbf{Skill area (internship exposure)} &
\textbf{How it is applied in this project} &
\textbf{Resulting deliverable} \\
\hline
Frontend development \newline (the internship work) &
Component-based UI, centralised state management, a real-time render loop, WebGL
rendering with Three.js, and a small HTTP/database service &
The implemented 3D game client and the leaderboard service (Chapter~4) \\
\hline
Business Analysis \newline (learned about, then practised here) &
Eliciting users and needs; decomposing the game into subsystems; writing
functional requirements, business rules and acceptance criteria; designing the
ranking policy &
The Software Requirements Specification and the game design (Chapter~3;
Appendix~A) \\
\hline
\end{tabular}
\end{table}

\section{Scope and Contributions}

This thesis covers a single, complete piece of work that we produced
independently. Its concrete contributions are:

\begin{itemize}
  \item \textbf{A complete Software Requirements Specification.} A structured
        catalogue of functional requirements, business rules, timing rules and
        non-functional requirements, each with acceptance criteria, covering
        every subsystem of the game (Appendix~A).
  \item \textbf{A game design derived from that specification.} The gameplay
        loop, the enemy and weapon balance, the area-of-effect damage model, the
        three-wave progression, and the leaderboard ranking rules (Chapter~3).
  \item \textbf{A working frontend application.} A real-time, interactive 3D
        game implemented with Next.js, React, TypeScript and Three.js/WebGL,
        organised into a clear layered architecture (Chapter~4).
  \item \textbf{A persistent leaderboard service.} A small server-side component
        --- HTTP endpoints backed by an SQLite database --- that records
        finished runs and returns them in ranked order (Chapter~4).
\end{itemize}

The scope is intentionally bounded. The game is a \emph{single-player},
\emph{desktop-browser} experience of \emph{three} waves, controlled with the
mouse. The following are explicitly out of scope: multiplayer or networked play,
user accounts and authentication, cross-device synchronisation, and mobile or
touch support. These boundaries keep the project achievable while still
exercising the full BA-to-implementation path, and several of them are revisited
as future work in Chapter~6.

\section{Report Structure}

The remainder of this report is organised as follows.

\begin{itemize}
  \item \textbf{Chapter~2 (Objectives)} states the general aim, the specific BA
        and Frontend objectives, and the success criteria used to judge the
        result.
  \item \textbf{Chapter~3 (Game Design)} presents the analysis work: the game
        concept, the requirements-engineering approach, the gameplay and
        economy design, the user-experience design, and the leaderboard ranking
        rules. It is the narrative companion to the full SRS in Appendix~A.
  \item \textbf{Chapter~4 (Materials and Methods)} describes the implementation:
        the technology stack and why each tool was chosen, the layered
        architecture and data flow, the state-management and game-loop design,
        the gameplay systems, the Three.js rendering layer, and the leaderboard
        service.
  \item \textbf{Chapter~5 (Results and Discussion)} evaluates the finished
        system against the specification through a requirements-traceability
        matrix, a gameplay walkthrough, and a performance assessment, and then
        reflects on the BA and Frontend decisions, the challenges met, and the
        project's limitations.
  \item \textbf{Chapter~6 (Conclusion and Perspective)} summarises what was
        achieved, names the BA and Frontend skills the project demonstrates, and
        outlines directions for future work.
  \item The \textbf{appendices} contain the full SRS, selected source-code
        listings, and a reference for the leaderboard API and database schema.
\end{itemize}
